\label{sec:limitations}

\subsection{Model Limitations}

\begin{enumerate}
    \item \textbf{Artificial Neural Network}: While transformers are useful models of neural computation, they differ from biological neural networks in architecture, learning rules, and dynamics. Our findings may not directly generalize to biological EEG.
    
    \item \textbf{Fixed Prompts}: We use a constant set of prompts to isolate the effect of sparsity. In realistic scenarios, content variation would interact with routing structure.
    
    \item \textbf{Linear Projection}: Real EEG involves nonlinear mixing through the head volume conductor. Future work should test whether nonlinear projections improve transfer.
    
    \item \textbf{No Temporal Dynamics}: We extract static pathway signatures. Real EEG is a time series; temporal correlations may further improve information transfer.
\end{enumerate}

\subsection{Experimental Design Limitations}

\begin{enumerate}
    \item \textbf{Single Model}: We only test GPT-2. Results may differ for larger models or different architectures (RNNs, CNNs, etc.).
    
    \item \textbf{Limited Sparsity Range}: Five discrete levels may not fully characterize the pathway-sparsity relationship. A continuous scan might reveal nonlinearities.
    
    \item \textbf{Small Sample Size}: 500 samples is modest by modern ML standards. Scaling to $10^4$--$10^5$ samples might reveal finer structure.
    
    \item \textbf{No Behavioral Data}: We do not measure whether changes in routing affect task performance. A triple dissociation (sparsity $\rightarrow$ routing $\rightarrow$ performance) would strengthen causal claims.
\end{enumerate}

\subsection{Statistical Limitations}

\begin{enumerate}
    \item \textbf{Correlation vs. Causation}: While our transfer experiment supports causality, we cannot rule out confounding. For example, both pathways and EEG might be downstream of a latent variable.
    
    \item \textbf{Generalization}: Validation is on held-out data from the same distribution. True generalization would require testing on novel prompts, models, or sparsity levels.
    
    \item \textbf{Multiple Comparisons}: We report 12+ correlations without correction for multiple comparisons. Bonferroni correction would reduce significance of some findings.
\end{enumerate}

\subsection{Future Directions}

\begin{enumerate}
    \item \textbf{Biological Validation}: Test predictions on real EEG data from humans performing cognit tasks. Can routing-based models explain EEG variance better than stimulus-based models?
    
    \item \textbf{Nonlinear Projections}: Develop learned nonlinear mappings from pathways to EEG, potentially improving transfer.
    
    \item \textbf{Temporal EEG}: Extend to time series by modeling dynamic attention and coupling to neural oscillations.
    
    \item \textbf{Other Modalities}: Test whether routing structure projects to MEG, fMRI, or electrocorticography (all have different spatial/temporal resolution).
    
    \item \textbf{Causal Manipulation}: Use TMS or pharmacological interventions to manipulate attention in humans and measure resulting EEG changes.
    
    \item \textbf{Larger Models}: Extend to GPT-3 scale or vision transformers to probe scaling behavior.
\end{enumerate}
